\documentclass[a4paper,oneside,bibtotoc,smallheadings,pointlessnumbers,
halfparskip,DIV15]{scrartcl}
% kate: remove-trailing-space on; replace-trailing-space-save on; indent-width 2; indent-mode normal; space-indent on;
% tocleft,
\usepackage[pdftex]{graphicx}
\usepackage[ngerman]{babel}
\usepackage[utf8]{inputenc}
\usepackage{../../assets/uniinput}
\usepackage[T1]{fontenc}
\usepackage{amssymb,amsmath}
\usepackage{booktabs}
\usepackage{enumerate}
\usepackage{subfig}
% \usepackage{siunitx}
% \sisetup{
%   seperr         = true,
%   trapambigerr   = true,
%   openerr        = (,
%   closeerr       = ),
%   expproduct     = cdot,
%   padnumber      = both,
%   stickyper      = true,
%   per            = reciprocal,
%   trapambigfrac  = true,
%   repeatunits    = false,
%   openfrac       = (,
%   closefrac      = ),
%   prefixsymbolic = true,
%   prefixproduct  = cdot,
%   decimalsymbol  = comma,
%   tabnumalign    = left,
%   tabtextalign   = left
% }
\usepackage{color}
\definecolor{grey}{rgb}{.4,.4,.4}
\definecolor{darkgreen}{rgb}{0,.35,0}
\definecolor{ltgray}{gray}{0.95}
% \definecolor{darkblue}{rgb}{0,0,.6}
% \definecolor{darkred}{rgb}{.6,0,0}
% \definecolor{red}{rgb}{.98,0,0}
\usepackage{listings}
\lstdefinelanguage{Maxima}{
  keywords={addrow,addcol,zeromatrix,ident,augcoefmatrix,ratsubst,diff,ev,tex,%
    with_stdout,nouns,express,depends,load,submatrix,div,grad,curl,%
    rootscontract,solve,part,assume,sqrt,integrate,abs,inf,exp,float,log},
  sensitive=true,
  comment=[n][\itshape]{/*}{*/}
}
\lstdefinelanguage{C++}{
%   commentstyle=\itshape\color{darkgreen},
%   commentstyle=\color{darkgreen},
%   keywordstyle=\bfseries, %\color{darkblue},
%   stringstyle=\color{darkred},
%   basicstyle=\ttfamily\scriptsize,
  morekeywords={TH1F,TLorentzVector,TVector3,vector,TFile,TFitResultPtr,TF1,\
                TGraph,TH1,TObject,TCanvas,string,Double_t,TGraphErrors},
  basicstyle=\scriptsize,
  numbers=left,
  numberstyle=\tiny,%\color{gray},
  stepnumber=1,
  tabsize=4,
  showspaces=false,
  showstringspaces=false,
  breaklines=true,
  frame=lrtb,
  captionpos=b,
  extendedchars=true,
  inputencoding=utf8,
%   backgroundcolor=\color{ltgray}
}
% \usepackage{courier}
\lstset{language=Matlab,
%     keywords={break,case,catch,continue,else,elseif,end,for,function,global,if,otherwise,persistent,return,switch,try,while},
%     basicstyle=\ttfamily\small,%\scriptsize,
  basicstyle=\scriptsize,
  keywordstyle=\bfseries,
  commentstyle=\itshape\color{blue},
  stringstyle=\itshape,
  identifierstyle=\ttfamily,
  numbers=left,
  numberstyle=\tiny,
  stepnumber=1,
  tabsize=2,
  showspaces=false,
  showstringspaces=false,
  breaklines=true,
  frame=single,
  columns=fixed,
  captionpos=b,
  extendedchars=\true,
  inputencoding=utf8,
  backgroundcolor=\color{ltgray},
  postbreak = \raisebox{0ex}[0ex][0ex]{\ensuremath{\hookrightarrow}}
%   linewidth=0.9\textwidth
}
\usepackage[pdftex]{hyperref}
\hypersetup{
%   colorlinks  = true,
%   urlcolor    = darkblue,
  pdftitle    = {Protokoll zum CPII-Übungsblatt 1},
  pdfsubject  = {Protokoll Computational Physics}
  pdfauthor   = {robert.riemann@physik.hu-berlin.de,thomas.murach@physik.hu-berlin.de},
  pdfkeywords = {Computational Physics,CP,II,2010,Jacobi-Methode,Diagonalisierung,Eigenwerte, Eigenvektoren,eig,octave,matlab}
%   pdfcreator  = {pdftex},
%   pdfproducer = {pdftex}
}

\newcommand{\dd}[1]{\mathrm{d}#1\,} % declare dx operator, usage: \dd{x} for dx
\newcommand{\lref}[1]{Listing (\ref{lst:#1})} % refer to a listing, usage: \lref{label} for Listing (...)
\newcommand{\fref}[1]{Abb. (\ref{fig:#1})} % refer to a figure, usage: \lref{label} for Abb. (...)
\newcommand{\tref}[1]{Tab. (\ref{tab:#1})} % refer to a table, usage: \lref{label} for Tab. (...)
\newcommand{\eref}[1]{Gl. (\ref{eqn:#1})} % refer to a equation, usage: \lref{label} for Tab. (...)

\begin{document}
% % % % % % % % % % % % % % % % % % % % % % % %
\title{{\centering \rule{15cm}{0.001cm}\\
\Large{\textsc{Institut für Physik der
Humboldt-Universität zu Berlin}}}\\ \centering \rule{15cm}{0.001cm}\\
\vspace{15mm} \centering
\includegraphics[scale=0.9]{../../assets/siegel}\\
\vspace{18mm}
{\bf{\huge{Computational Physics II}}}\\
\vspace{12mm}
Übungsblatt 3\\
\vspace{15mm}
% Compton-Effekt\\
% \vspace{14mm} {\small{\textbf{Betreuer: M. zur Nedden}}}\\}
}
\author{Robert Riemann; Matr.Nr.: 521085\\
Thomas Murach; Matr.Nr.: 517771\vspace{18mm}}
% \vspace{8mm}
% \date{15. Juni 2008}
% % % % % % % % % % % % % % % % % % % % % % % %
% \onecolumn
\maketitle
% \twocolumn

\newpage
% \tableofcontents
% \listoffigures
% \listoftables
\section*{Aufgabe 3.1a}
Im ersten Teil der Aufgabe war zu untersuchen, wie sich die logistische
Abbildung $$M(x) = rx\cdot (1-x)$$ für $r$-Werte von 4, 3.9 und 3.832 verhält.
Im Speziellen wurde das invariante Maß $ϱ(x)$ untersucht, indem für 1000 zu
iterierende Punkte die erreichten $x$-Werte histogrammiert wurden. Der
Quelltext, der dies umsetzt, ist in \lref{logabb} dargestellt.

\lstinputlisting[firstline=1,firstnumber=1,label=lst:logabb,caption={logabb.m}]{
../logabb.m}

Die hierin aufgerufenen Funktionen sind in \lref{logmap} und \lref{logmapp}
dargestellt.

\lstinputlisting[firstline=1,firstnumber=1,label=lst:logmap,caption={logmap.m}]{
../logmap.m}

\lstinputlisting[firstline=1,firstnumber=1,label=lst:logmapp,caption={logmapp.m}
]{../logmapp.m}

Die Vorlagen für die Quelltexte stammen aus der Vorlesung. 

Die resultierenden Plots sind in \fref{logabb} dargestellt.
\begin{figure}[htb]
\centering
  \includegraphics[width=1\columnwidth,keepaspectratio]{../logabb.png}
  \caption{Output des Programms}
  \label{fig:logabb}
\end{figure}
% Nun sollen die in der Aufgabenstellung gegebenen Relationen
\begin{eqnarray}
v'_{ip} &=& v_{ip} - s(v_{iq} + τv_{ip})\\
v'_{iq} &=& v_{iq} + s(v_{ip} - τv_{iq})
\end{eqnarray}
bewiesen werden, wobei $v_{ab}$ bzw. $v'_{ab}$ die $(a,b)$-Komponente der
Matrizen $V$ bzw. $V'$ sind. Um diese Relationen zu erhalten, muss der
konkrete Aufbau der Matrizen $P$ bekannt sein. Diese bestehen aus einer
Einheitsmatrix, bei der einzig die Elemente in der $p$- bzw. $q$-ten Zeile bzw.
Spalte verändert werden. Dies lässt sich wie folgt darstellen:
\begin{equation}
p_{ij} = \begin{cases}\begin{array}[h]{ll}
1 & \mathrm{f"ur~} i=j \notin \{p,q\} \\c & \mathrm{f"ur~} i=j \in \{p,q\}\\s &
\mathrm{f"ur~} i=p, j=q\\-s & \mathrm{f"ur~} i=q, j=p\\0 & \mathrm{sonst}
\end{array}\end{cases}
\end{equation}
% auch interessant: (aber nur außerhalb der mathe-umgebung
% \begin{tabbing}
% 1 \= \quad \= iae\\c \> \> iaee\\s \> \> aie\\-s \> \> ie\\0 \> \> ieie
% \end{tabbing}

In dieser Gleichung für $p_{ij}$ stehen $c$ bzw. $s$ für den Kosinus bzw. Sinus
des Drehwinkels φ:
\begin{eqnarray}
c &=& \cos(φ)\\
s &=& \sin(φ)
\end{eqnarray}

Um die Definition der $P$-Matrizen in eine kompaktere Form zu bringen, wird nun
eine Indexschreibweise mit Hilfe der Kronecker-Deltas $δ_{ij}$ verwendet:
\begin{eqnarray}
p_{ij} &=& δ_{ij} + s\cdot δ_{ip}δ_{jq} - s\cdot δ_{iq}δ_{jp} + (c-1)\cdot
δ_{ip}δ_{jp} + (c-1)\cdot δ_{iq}δ_{jq}
\end{eqnarray}
Von den beiden Diagonalelementen muss jeweils eine eins abgezogen werden, um
die durch den $δ_{ij}$-Term zu kompensieren. Wenn man nun eine „beliebige“
Matrix $V$ mit $P$ multipliziert, ergeben sich (unter Verwendung der
Einsteinschen Summenkonvention) die gesuchten Relationen:
\begin{eqnarray}
v'_{kj} &=& v_{ki} P_{ij} = v_{ki}δ_{ij} + s\cdot v_{ki}δ_{ip}δ_{jq} - s\cdot
v_{ki}δ_{iq}δ_{jp} + (c-1)v_{ki}(δ_{ip}δ_{jp} + δ_{iq}δ_{jq})\\
&=& v_{kj} + s\cdot v_{kp}δ_{jq} - s\cdot v_{kq}δ_{jp} + (c-1)(v_{kp}δ_{jp} +
v_{kq}δ_{jq})
\end{eqnarray}

Nun kann eine Fallunterscheidung durchgeführt werden, die berücksichtigt,
welche der Kronecker-Deltas null bzw. eins werden.

\begin{eqnarray}
\begin{array}[h]{cclcccl}
j &\notin& {p,q} &:& v'_{kj} &=& v_{kj}\\
j &=& p &:& v'_{kp} &=& v_{kp} - s\cdot v_{kq} + (c-1)v_{kp}\\
&&&&&=& v_{kp} - s\cdot(v_{kq} - \frac{c-1}{s}v_{kp})\\
j &=& q &:& v'_{kq} &=& v_{kq} + s\cdot v_{kp} + (c-1)v_{kq}\\
&&&&&=& v_{kp} + s\cdot(v_{kq} + \frac{c-1}{s}v_{kp})
\end{array}
\label{eqn:vmatrix}
\end{eqnarray}

Schließlich wird gezeigt, dass man den jeweils auftretenden Bruch in die
gewünschte Form bringen kann:
\begin{eqnarray}
\frac{c-1}{s} &=& \frac{s\cdot(c-1)}{s^2} = \frac{s\cdot(c-1)}{1-c^2} =
\frac{s\cdot(c-1)}{(1-c)\cdot(1+c)}\\
&=& -\frac{s}{c+1} \equiv -τ 
\end{eqnarray}

Dieses Ergebnis wird nun in \eref{vmatrix} eingesetzt:
\begin{eqnarray}
v'_{kp} &=& v_{kp} - s\cdot(v_{kq} + τv_{kp})\\
v'_{kq} &=& v_{kp} + s\cdot(v_{kq} - τv_{kp})
\end{eqnarray}
Dies entspricht der gesuchten Lösung.
% \section*{Aufgabe 2.1d)}
Schließlich sollten die systematischen Fehler abgeschätzt werden, indem man zum
einen ε variiert und zum anderen statt \texttt{ode45} \texttt{ode23} verwendet,
was prinzipiell ungenauere Ergebnisse liefert, oder auch direkt für
\texttt{ode45} selbst andere Genauigkeiten fordert.

\subsection*{Variation von ε}
Wenn man die Breite der normalverteilten Zufallszahlen, die die Abweichung
repräsentieren, ändert, so ergeben sich auch andere Zeitpunkte, bei denen im 
Mittel die Abweichung der Trajektorien größer als 0,1 wird. Hier wurde ε so
variiert, dass es sehr klein gegenüber den Werten von $y$ und $\dot{y}$ ist und
deutlich größer als die Maschinengenauigkeit. Eine genaue Analyse des Einflusses 
der Variation ist in \lref{epsilon} zu finden.

\lstinputlisting[firstline=1,firstnumber=1,label=lst:epsilon,caption={Auswirkungen 
der Wahl von ε auf die interessanten Größen}]{../syst_errors.dat}

Wie man erkennen kann, nimmt die Zeit, die benötigt wird, um die gewünschte 
Abweichung herbeizuführen, für kleiner werdendes ε stetig zu, was auch zu
erwarten ist, da bei kleinerem ε die Ausgangspunkte und -geschwindigkeiten der
bieden Trajektorien immer weniger voneinander abweichen. Interessant zu beobachten
ist die Tatsache, dass \texttt{t\_mean} für jedes nächstkleinere ε um etwa 20
Einheiten zunimmt. Dies könnte in einer tiefergehenden Analyse begründet werden.

Weiterhin werden die 
Verteilungen immer breiter, was sich in größeren Standardabweichungen und
Vertrauensintervallen niederschlägt. Das bedeutet, dass bei die genaue Wahl der
Abweichung in Bezug auf Richtung und Betrag für kleiner werdendes ε größeren
Einfluss gewinnt.

Schon die Wahl der nächstgrößeren bzw. -kleineren Größe von ε liefert deutlich
andere Ergebnisse für \texttt{t\_mean}. Die Abweichungen betragen jeweils ca.
17 Prozent. 

Die Werte von λ schwanken von 0,1284 und 0,1360 und stimmen somit innerhalb der
Standardabweichungen, die zwischen 0.025 und 0,07 liegen, näherungsweise überein.
Die Wahl von ε hat hier also nur einen geringen Einfluss. Einzig die
Standardabweichungen selbst werden für abnehmendes ε immer kleiner.

Weiterhin ist der Einfluss der konkreten Wahl der Implementierung der numerischen
Lösung der Differentialgleichung untersucht worden, indem statt \texttt{ode45}
die Funktion \texttt{ode23} verwendet wurde. Hierfür wurden wieder 
Standardparameter verwendet. Wie man erkennen kann, stimmen die Ergebnisse bis
auf ca. 1 Prozent überein, und die Standardabweichungen nehmen ebenfalls beinah
gleiche Werte an. Die Wahl der Funktion liefert demnach nur Unterschiede im 
Prozentbereich und kann somit gegenüber dem Einfluss der Wahl von ε
vernachlässigt werden.

Schließlich wurden andere relative und absolute Genauigkeiten gefordert. Hierfür
wurden wiederum Standardparameter gewählt (\texttt{ode45}, $ε=10^{-6}$). Dies
führt den am Ende in \lref{epsilon} dargestellten Ergebnissen.

Man kann erkennen, dass man durchaus eine Größenordnung von der Standard-Toleranz
abweichen kann, ohne dass dies größere Auswirkungen auf \texttt{t\_mean} etc.
hätte. Nur wenn man die Toleranz zu groß wählt resultieren auch stärkere
Abweichungen in der Größenordnung von fünf Prozent.

Der dominante systematische Fehler ist also in in der konkreten Wahl von ε zu
suchen. Dieser beläuft sich auf ca. 17 \%, wenn man nur einen Bereich von $ε = 
10^{-5}$ bis $ε = 10^{-7}$ zulässt.

Die durch die verschiedenen Untersuchungen erhaltenen λ-Werte variieren zwischen
0,127 ± 0,042 und 0,13599 ± 0,068 und sind somit im gesamten Bereich der
Veränderungen der Parameter näherungsweise konstant geblieben.

Eine weitere Erhöhung der Anzahl der Wiederholungen würde zu einer Erhöhung der 
statistischen Genauigkeit führen. Da die statistischen Fehler, repräsentiert
durch das Vertrauensintervall, im Bereich von ca. zwei Prozent liegt, dominiert
insgesamt der systematische Fehler. Demnach ist es nicht sinnvoll, noch mehr 
Versuche durchzuführen, um den gesamten Fehler deutlich zu senken.
% \input{auswertung}
% \input{ergebnis}
% \input{quellen}
% \input{anhang}

% \section{Beispiele}
% blabla
%
% Weitere Information sind der Versuchsbeschreibung im \cite{script} zu entnehmen.
%
% \lstinputlisting[firstline=2,firstnumber=1,label=lst:aufruf,caption={aufruf.m - Matlab Batchdatei}]{aufruf.m}
%
% \begin{eqnarray}
% 	g &=& \SI{9.8157(11)}{\metre\per\second\squared} \\
%         \left[ g \right] &=& \si{\metre\per\Square\second} % großschreibung!
% 	\label{eqn:asdf}
% \end{eqnarray}
%
% Dann kann man auf \eref{asdf} verweisen.


% \begin{figure}[htb]
% 	\centering
% 	\includegraphics[width=1\columnwidth,keepaspectratio]{Winkelabhaengigkeit2}
% 	\caption{Periodendauer in Abhängigkeit der Maximalauslenkung}
% 	\label{fig:Winkelabhaengigkeit}
% \end{figure}

% \begin{table}[htbp]
% \centering
% \setlength{\tabcolsep}{14pt}
% \begin{tabular*}{\columnwidth}{%
% S[tabformat=2.1]%
% S[tabformat=2.2]%
% S[tabformat=1.2]}
% \toprule
% {$R$ in \si{\ohm}} &
% {$U$ in \si{\volt}} &
% {$\frac{U}{U_{R=0}}$}\\
% \midrule
% \multicolumn{3}{c}{\textit{Eingangswiderstand $R_E$}}\\
% \midrule
% 0 & 16 & 1 \\
% 4.7e3 & 8 & 0.5 \\
% \midrule
% \multicolumn{3}{c}{\textit{Ausgangswiderstand $R_A$}}\\
% \midrule
% 0 & 1.8 & 1 \\
% 12 & 0.52 & 0.29 \\
% 27 & 1.2 & 0.67 \\
% 18 & 0.8 & 0.44 \\
% \bottomrule
% \end{tabular*}
% \label{tab:221messungexp_widerstaende}
% \caption{experimentelle Messung der Ein- und Ausgangswiderstände}
% \end{table}


% \newpage

% \onecolumn
% \appendix
% \begin{figure}
% 	\centering
% 	\includegraphics[width=0.98\textwidth,keepaspectratio]{Messprotokoll}
% 	\caption{Messprotokoll}
% 	\label{fig:protokoll}
% \end{figure}
%\colorbox{yellow}{} Farben verwenden
\end{document}
